\section{Results and Analysis}
\label{sec:results}

\paragraph{Execution summary and reporting contract.}

All 32 cells completed on eight \HevHardware{} in FP32 and used
\HevGPUHours{} GPU-hours. Table~\ref{tab:dev-matrix} reports corrected
development-selection F1. Each cell selects its checkpoint and operating
threshold on the same eight-scene subset. These scores therefore include
model-selection optimism and are not held-out estimates.
Figure~\ref{fig:label-efficiency} shows the four-point curves, and
Fig.~\ref{fig:training-dynamics} shows their training histories in the
appendix. The evidence generator reconstructs every value from the frozen
cohort, checkpoint-bound completion records, and full training curves. It
rejects an incomplete cohort or a value that disagrees with its source
record.

\begin{table}[t]
\centering
\caption{Development-selection F1 for all 32 cells at each cell's
checkpoint-bound operating threshold. Budgets are fractions of the 111
training scenes (scene counts in parentheses). Per-cell thresholds, best
epochs, and held-out slots appear in Appendix~\ref{app:inventory}.}
\label{tab:dev-matrix}
\small
\setlength{\tabcolsep}{5pt}
\begin{tabular}{@{}llcccc@{}}
\toprule
Track & Initialization & 10\% (12) & 25\% (28) & 50\% (56) & 100\% (111) \\
\midrule
ViT-B/16 & Random & \HevDevFViTRandomTen & \HevDevFViTRandomTwentyFive & \HevDevFViTRandomFifty & \HevDevFViTRandomHundred \\
 & Optical & \HevDevFViTOpticalTen & \HevDevFViTOpticalTwentyFive & \HevDevFViTOpticalFifty & \HevDevFViTOpticalHundred \\
 & SAR & \HevDevFViTSarTen & \HevDevFViTSarTwentyFive & \HevDevFViTSarFifty & \HevDevFViTSarHundred \\
 & ImageNet & \HevDevFViTImageNetTen & \HevDevFViTImageNetTwentyFive & \HevDevFViTImageNetFifty & \HevDevFViTImageNetHundred \\
\midrule
ConvNeXt-V2-B & Random & \HevDevFCNNRandomTen & \HevDevFCNNRandomTwentyFive & \HevDevFCNNRandomFifty & \HevDevFCNNRandomHundred \\
 & Optical & \HevDevFCNNOpticalTen & \HevDevFCNNOpticalTwentyFive & \HevDevFCNNOpticalFifty & \HevDevFCNNOpticalHundred \\
 & SAR & \HevDevFCNNSarTen & \HevDevFCNNSarTwentyFive & \HevDevFCNNSarFifty & \HevDevFCNNSarHundred \\
 & ImageNet & \HevDevFCNNImageNetTen & \HevDevFCNNImageNetTwentyFive & \HevDevFCNNImageNetFifty & \HevDevFCNNImageNetHundred \\
\bottomrule
\end{tabular}
\end{table}

\begin{figure}[t]
\centering
\includegraphics[width=\textwidth]{heldout_label_efficiency.pdf}
\caption{Development-F1 label-efficiency curves, one panel per track. Solid
lines are pretrained arms. The dashed gray line is the track's random floor.
Values are seed-0 point estimates. We draw no uncertainty ribbon.}
\label{fig:label-efficiency}
\end{figure}

\paragraph{Label-efficiency curves and transfer gains.}

At 12 scenes, all three pretrained ViTs beat the random ViT. Their gains are
\HevDeltaFTenViTOptical{} for optical, \HevDeltaFTenViTSar{} for SAR, and
\HevDeltaFTenViTImageNet{} for ImageNet (Eq.~\eqref{eq:gain} and
Fig.~\ref{fig:transfer-gains}). In the CNN track, SAR gains
\HevDeltaFTenCNNSar{} and ImageNet gains \HevDeltaFTenCNNImageNet{} over the
12-scene random floor. The optical CNN matches random at 12 scenes
(both \HevDevFCNNOpticalTen) and trails it at every larger budget. Transfer
benefit is therefore specific to the checkpoint rather than guaranteed by
pretraining.

More labeled scenes did not raise every development curve. ViT-optical falls
from \HevDevFViTOpticalTen{} to \HevDevFViTOpticalTwentyFive{} between the
first two budgets. ViT-SAR falls from \HevDevFViTSarTen{} to
\HevDevFViTSarTwentyFive. We report these declines without smoothing or
reruns. With one seed and 517 positives on the selection subset, this study
cannot separate optimization variation, subset composition, and selection
error.

We also measure the smallest tested budget that reaches the track's
full-data random floor. CNN-ImageNet reaches that reference at 28 scenes
(\HevDevFCNNImageNetTwentyFive{} versus
\HevDevFCNNRandomHundred). Every ViT transfer arm and CNN-SAR reaches it at
56 scenes. CNN-optical never does. At 12 scenes, the strongest transferred
model in each track falls just short of its full-data random floor. The
tested checkpoints can therefore approach, but do not match, the selected
development score of a random model trained on all 111 scenes.

\paragraph{Source-domain contrasts and architecture generality.}

The SAR--optical ranking changes with the backbone
(Eq.~\eqref{eq:contrast} and Fig.~\ref{fig:transfer-gains}). CNN-SAR leads
CNN-optical at every budget. The largest gap occurs at 12 scenes
(\HevDevFCNNSarTen{} versus \HevDevFCNNOpticalTen). The BigEarthNet
Sentinel-1 and Sentinel-2 checkpoints share their architecture, dataset
family, and supervised task, which makes this the cleanest modality
comparison in the study.

The ViT ranking reverses below half data. Optical leads SAR at 10\% and 25\%
(\HevDevFViTOpticalTen{} versus \HevDevFViTSarTen, and
\HevDevFViTOpticalTwentyFive{} versus \HevDevFViTSarTwentyFive). SAR moves
ahead at 56 and 111 scenes. The scarce-label SAR--optical order therefore
does not replicate across architectures. The ViT checkpoints also differ in
objective and source corpus, so their ranking cannot identify a modality-only
effect.

At 111 scenes, development selection ranks ImageNet first in both tracks
(\HevDevFViTImageNetHundred{} and \HevDevFCNNImageNetHundred). Held-out
scoring tests whether these full-data rankings persist.

\paragraph{Held-out results.}

\ifHevTestComplete
All 32 cells were scored once on the frozen 16-scene test split with their
development-bound thresholds and \HevTestInferencePrecision{} inference.
The test column appears only after every result validates against the frozen
cohort. Test F1 is lower than development F1 for every cell. The gap ranges
from 0.041 to 0.142 and averages \HevTestMeanGap{} F1. A negative gap is
compatible with selection on eight scenes, but this experiment cannot
separate selection optimism, threshold transfer, and scene-distribution differences.

The central SAR--optical ranking persists at every budget. CNN-SAR remains
above CNN-optical throughout. In ViT, optical remains ahead at 10\% and 25\%,
then SAR leads at 50\% and 100\%. The full-data ranking changes in one track.
ViT-SAR reaches \HevTestFViTSarHundred, above ImageNet's
\HevTestFViTImageNetHundred. CNN-ImageNet retains its lead at
\HevTestFCNNImageNetHundred.

The predeclared monotonicity gate failed on two SAR test curves. ViT-SAR
falls from \HevTestFViTSarTen{} at 10\% to
\HevTestFViTSarTwentyFive{} at 25\%. CNN-SAR rises on development from
\HevDevFCNNSarFifty{} at 50\% to \HevDevFCNNSarHundred{} at 100\%, but its
test F1 falls from \HevTestFCNNSarFifty{} to
\HevTestFCNNSarHundred. Both test declines exceed the allowed 0.02. The
experiment therefore does not establish a monotone label-efficiency curve.
Appendix~\ref{app:inventory} lists the full test column.
\else
The 16-scene test split has not been scored in this build. The held-out
contract first freezes the 32-cell cohort. It then applies each cell's
development-bound threshold and writes one immutable result. The generator
admits test values only when all 32 results validate.
\fi
\ifHevFinalEval
The sealed 50-scene human-verified evaluation ran once for all 32 cells.
Dark-vessel recall and near-shore results appear in
Appendix~\ref{app:inventory}.
\else
The 50 human-verified scenes remain sealed. Dark-vessel recall and near-shore
performance are defined only on that set, so neither result is available in
this report.
\fi
